% Accenture x Faculty co-brand marks (concept 01 "chevron crossbar").
% Original vector drawings for internal review. Not official artwork of either
% firm; any external use needs sign-off from both brand teams.
%
% Build one file:  pdflatex -jobname=<name> "\def\V{A}\def\S{dark}% Accenture x Faculty co-brand marks (concept 01 "chevron crossbar").
% Original vector drawings for internal review. Not official artwork of either
% firm; any external use needs sign-off from both brand teams.
%
% Build one file:  pdflatex -jobname=<name> "\def\V{A}\def\S{dark}% Accenture x Faculty co-brand marks (concept 01 "chevron crossbar").
% Original vector drawings for internal review. Not official artwork of either
% firm; any external use needs sign-off from both brand teams.
%
% Build one file:  pdflatex -jobname=<name> "\def\V{A}\def\S{dark}% Accenture x Faculty co-brand marks (concept 01 "chevron crossbar").
% Original vector drawings for internal review. Not official artwork of either
% firm; any external use needs sign-off from both brand teams.
%
% Build one file:  pdflatex -jobname=<name> "\def\V{A}\def\S{dark}\input{logos.tex}"
%   \V = A (master lockup), D (integral stem), E (risk to trust), I (co-brand lockup)
%   \S = dark | light | white | black   (white/black = transparent background)
% build.sh builds every combination and exports PDF, SVG and PNG.
\documentclass[tikz,border=0pt]{standalone}
\usepackage[T1]{fontenc}
\usepackage{lmodern}
\usepackage{amsmath}
\usepackage[letterspace=120]{microtype}
\usetikzlibrary{calc,backgrounds}
\renewcommand{\familydefault}{\sfdefault}

\definecolor{acnPurple}{HTML}{A100FF}
\definecolor{acnNight}{HTML}{09090B}
\definecolor{acnGrey}{HTML}{71717A}

% House ">" (same geometry as the decks), bottom-left #1, height #2
\def\chevpath#1#2{($#1+(0,#2)$) -- ($#1+(0.62*#2,0.5*#2)$) -- #1 -- ($#1+(0,0.205*#2)$)
   -- ($#1+(0.365*#2,0.5*#2)$) -- ($#1+(0,0.795*#2)$) -- cycle}
% The same ">" grown by a 4.5% keyline (exact miter offset), clockwise
\def\chevgap#1#2{($#1+(-0.0450*#2,1.0941*#2)$) -- ($#1+(0.6917*#2,0.5000*#2)$)
   -- ($#1+(-0.0450*#2,-0.0941*#2)$) -- ($#1+(-0.0450*#2,0.2265*#2)$)
   -- ($#1+(0.2934*#2,0.5000*#2)$) -- ($#1+(-0.0450*#2,0.7735*#2)$) -- cycle}
% Clip to box #1--#2 minus the keyline area (a real gap, so transparent versions work)
\newcommand{\gapclip}[4]{\clip #3 rectangle #4 \chevgap{#1}{#2};}

% Tapered S stem (Faculty-style, redrawn). Unit box 2 x 1.5.
\def\stempath{(0,-0.018) -- (0.5,-0.03) .. controls (1.045,-0.03) and (1.045,1.47) .. (1.5,1.47)
   -- (2,1.482) -- (2,1.518) -- (1.5,1.53) .. controls (0.955,1.53) and (0.955,0.03) .. (0.5,0.03)
   -- (0,0.018) -- cycle}

\newcommand{\facword}[2][11]{{\fontsize{#1}{#1}\selectfont\color{#2}\textls[260]{FACULTY}}}
\newcommand{\acnword}[2][14]{{\fontsize{#1}{#1}\selectfont\bfseries\color{#2}accenture}}

% Integral stem with chevron crossbar. #1 centre, #2 glyph scale, #3 chevron height, #4 fg
\newcommand{\intmark}[4]{%
  \begin{scope}\gapclip{($#1+(-0.24*#3,-0.5*#3)$)}{#3}{($#1+(-0.25*#2,-0.32*#2)$)}{($#1+(0.25*#2,0.32*#2)$)}
    \node[inner sep=0,text=#4,scale=#2] at #1 {$\displaystyle\int$};
  \end{scope}
  \fill[acnPurple] \chevpath{($#1+(-0.24*#3,-0.5*#3)$)}{#3};}

% ---- the four marks, #1 = foreground colour ----
\newcommand{\markA}[1]{% 01A master lockup
  \begin{scope}\gapclip{(0.6715,0.3)}{0.9}{(0,-0.05)}{(2,1.55)}\fill[#1] \stempath;\end{scope}
  \fill[acnPurple] \chevpath{(0.6715,0.3)}{0.9};
  \node[anchor=base,inner sep=0] at (1.0,-0.83) {\facword{#1}};}
\newcommand{\markD}[1]{% 01D integral stem
  \intmark{(0,0)}{6.2}{1.0}{#1}}
\newcommand{\markE}[1]{% 01E from risk to trust
  \intmark{(0,0)}{6.2}{1.0}{#1}
  \node[anchor=west,inner sep=0,text=#1,scale=1.5] at (0.67,1.35) {$\mathit{trust}$};
  \node[anchor=west,inner sep=0,text=#1,scale=1.5] at (-0.03,-1.48) {$\mathit{risk}$};
  \node[anchor=west,inner sep=0,text=acnGrey,scale=2.2] at (0.95,-0.03) {$\mathrm{d}\mathit{AI}$};}
\newcommand{\markI}[1]{% 01I co-brand lockup
  \node[anchor=base east,inner sep=0] (a) at (-0.75,-0.22) {\acnword{#1}};
  \fill[acnPurple] \chevpath{($(a.base east)+(-0.36,0.43)$)}{0.2};
  \draw[#1,line width=0.5pt] (-0.48,-0.65) -- (-0.48,0.65);
  \intmark{(0,0)}{2.9}{0.48}{#1}
  \node[anchor=base west,inner sep=0] at (0.45,-0.22) {\facword[10]{#1}};}

\def\fgcol{white}\def\bgcol{}
\def\tmp{light}\ifx\S\tmp\def\fgcol{black}\def\bgcol{white}\fi
\def\tmp{dark}\ifx\S\tmp\def\fgcol{white}\def\bgcol{acnNight}\fi
\def\tmp{black}\ifx\S\tmp\def\fgcol{black}\fi

\begin{document}
\begin{tikzpicture}[x=14mm,y=14mm]
  \begin{scope}[local bounding box=mark]
    \csname mark\V\endcsname{\fgcol}
  \end{scope}
  \ifx\bgcol\empty
    \useasboundingbox ($(mark.south west)+(-1mm,-1mm)$) rectangle ($(mark.north east)+(1mm,1mm)$);
  \else
    \begin{scope}[on background layer]
      \fill[\bgcol] ($(mark.south west)+(-9mm,-9mm)$) rectangle ($(mark.north east)+(9mm,9mm)$);
    \end{scope}
  \fi
\end{tikzpicture}
\end{document}
"
%   \V = A (master lockup), D (integral stem), E (risk to trust), I (co-brand lockup)
%   \S = dark | light | white | black   (white/black = transparent background)
% build.sh builds every combination and exports PDF, SVG and PNG.
\documentclass[tikz,border=0pt]{standalone}
\usepackage[T1]{fontenc}
\usepackage{lmodern}
\usepackage{amsmath}
\usepackage[letterspace=120]{microtype}
\usetikzlibrary{calc,backgrounds}
\renewcommand{\familydefault}{\sfdefault}

\definecolor{acnPurple}{HTML}{A100FF}
\definecolor{acnNight}{HTML}{09090B}
\definecolor{acnGrey}{HTML}{71717A}

% House ">" (same geometry as the decks), bottom-left #1, height #2
\def\chevpath#1#2{($#1+(0,#2)$) -- ($#1+(0.62*#2,0.5*#2)$) -- #1 -- ($#1+(0,0.205*#2)$)
   -- ($#1+(0.365*#2,0.5*#2)$) -- ($#1+(0,0.795*#2)$) -- cycle}
% The same ">" grown by a 4.5% keyline (exact miter offset), clockwise
\def\chevgap#1#2{($#1+(-0.0450*#2,1.0941*#2)$) -- ($#1+(0.6917*#2,0.5000*#2)$)
   -- ($#1+(-0.0450*#2,-0.0941*#2)$) -- ($#1+(-0.0450*#2,0.2265*#2)$)
   -- ($#1+(0.2934*#2,0.5000*#2)$) -- ($#1+(-0.0450*#2,0.7735*#2)$) -- cycle}
% Clip to box #1--#2 minus the keyline area (a real gap, so transparent versions work)
\newcommand{\gapclip}[4]{\clip #3 rectangle #4 \chevgap{#1}{#2};}

% Tapered S stem (Faculty-style, redrawn). Unit box 2 x 1.5.
\def\stempath{(0,-0.018) -- (0.5,-0.03) .. controls (1.045,-0.03) and (1.045,1.47) .. (1.5,1.47)
   -- (2,1.482) -- (2,1.518) -- (1.5,1.53) .. controls (0.955,1.53) and (0.955,0.03) .. (0.5,0.03)
   -- (0,0.018) -- cycle}

\newcommand{\facword}[2][11]{{\fontsize{#1}{#1}\selectfont\color{#2}\textls[260]{FACULTY}}}
\newcommand{\acnword}[2][14]{{\fontsize{#1}{#1}\selectfont\bfseries\color{#2}accenture}}

% Integral stem with chevron crossbar. #1 centre, #2 glyph scale, #3 chevron height, #4 fg
\newcommand{\intmark}[4]{%
  \begin{scope}\gapclip{($#1+(-0.24*#3,-0.5*#3)$)}{#3}{($#1+(-0.25*#2,-0.32*#2)$)}{($#1+(0.25*#2,0.32*#2)$)}
    \node[inner sep=0,text=#4,scale=#2] at #1 {$\displaystyle\int$};
  \end{scope}
  \fill[acnPurple] \chevpath{($#1+(-0.24*#3,-0.5*#3)$)}{#3};}

% ---- the four marks, #1 = foreground colour ----
\newcommand{\markA}[1]{% 01A master lockup
  \begin{scope}\gapclip{(0.6715,0.3)}{0.9}{(0,-0.05)}{(2,1.55)}\fill[#1] \stempath;\end{scope}
  \fill[acnPurple] \chevpath{(0.6715,0.3)}{0.9};
  \node[anchor=base,inner sep=0] at (1.0,-0.83) {\facword{#1}};}
\newcommand{\markD}[1]{% 01D integral stem
  \intmark{(0,0)}{6.2}{1.0}{#1}}
\newcommand{\markE}[1]{% 01E from risk to trust
  \intmark{(0,0)}{6.2}{1.0}{#1}
  \node[anchor=west,inner sep=0,text=#1,scale=1.5] at (0.67,1.35) {$\mathit{trust}$};
  \node[anchor=west,inner sep=0,text=#1,scale=1.5] at (-0.03,-1.48) {$\mathit{risk}$};
  \node[anchor=west,inner sep=0,text=acnGrey,scale=2.2] at (0.95,-0.03) {$\mathrm{d}\mathit{AI}$};}
\newcommand{\markI}[1]{% 01I co-brand lockup
  \node[anchor=base east,inner sep=0] (a) at (-0.75,-0.22) {\acnword{#1}};
  \fill[acnPurple] \chevpath{($(a.base east)+(-0.36,0.43)$)}{0.2};
  \draw[#1,line width=0.5pt] (-0.48,-0.65) -- (-0.48,0.65);
  \intmark{(0,0)}{2.9}{0.48}{#1}
  \node[anchor=base west,inner sep=0] at (0.45,-0.22) {\facword[10]{#1}};}

\def\fgcol{white}\def\bgcol{}
\def\tmp{light}\ifx\S\tmp\def\fgcol{black}\def\bgcol{white}\fi
\def\tmp{dark}\ifx\S\tmp\def\fgcol{white}\def\bgcol{acnNight}\fi
\def\tmp{black}\ifx\S\tmp\def\fgcol{black}\fi

\begin{document}
\begin{tikzpicture}[x=14mm,y=14mm]
  \begin{scope}[local bounding box=mark]
    \csname mark\V\endcsname{\fgcol}
  \end{scope}
  \ifx\bgcol\empty
    \useasboundingbox ($(mark.south west)+(-1mm,-1mm)$) rectangle ($(mark.north east)+(1mm,1mm)$);
  \else
    \begin{scope}[on background layer]
      \fill[\bgcol] ($(mark.south west)+(-9mm,-9mm)$) rectangle ($(mark.north east)+(9mm,9mm)$);
    \end{scope}
  \fi
\end{tikzpicture}
\end{document}
"
%   \V = A (master lockup), D (integral stem), E (risk to trust), I (co-brand lockup)
%   \S = dark | light | white | black   (white/black = transparent background)
% build.sh builds every combination and exports PDF, SVG and PNG.
\documentclass[tikz,border=0pt]{standalone}
\usepackage[T1]{fontenc}
\usepackage{lmodern}
\usepackage{amsmath}
\usepackage[letterspace=120]{microtype}
\usetikzlibrary{calc,backgrounds}
\renewcommand{\familydefault}{\sfdefault}

\definecolor{acnPurple}{HTML}{A100FF}
\definecolor{acnNight}{HTML}{09090B}
\definecolor{acnGrey}{HTML}{71717A}

% House ">" (same geometry as the decks), bottom-left #1, height #2
\def\chevpath#1#2{($#1+(0,#2)$) -- ($#1+(0.62*#2,0.5*#2)$) -- #1 -- ($#1+(0,0.205*#2)$)
   -- ($#1+(0.365*#2,0.5*#2)$) -- ($#1+(0,0.795*#2)$) -- cycle}
% The same ">" grown by a 4.5% keyline (exact miter offset), clockwise
\def\chevgap#1#2{($#1+(-0.0450*#2,1.0941*#2)$) -- ($#1+(0.6917*#2,0.5000*#2)$)
   -- ($#1+(-0.0450*#2,-0.0941*#2)$) -- ($#1+(-0.0450*#2,0.2265*#2)$)
   -- ($#1+(0.2934*#2,0.5000*#2)$) -- ($#1+(-0.0450*#2,0.7735*#2)$) -- cycle}
% Clip to box #1--#2 minus the keyline area (a real gap, so transparent versions work)
\newcommand{\gapclip}[4]{\clip #3 rectangle #4 \chevgap{#1}{#2};}

% Tapered S stem (Faculty-style, redrawn). Unit box 2 x 1.5.
\def\stempath{(0,-0.018) -- (0.5,-0.03) .. controls (1.045,-0.03) and (1.045,1.47) .. (1.5,1.47)
   -- (2,1.482) -- (2,1.518) -- (1.5,1.53) .. controls (0.955,1.53) and (0.955,0.03) .. (0.5,0.03)
   -- (0,0.018) -- cycle}

\newcommand{\facword}[2][11]{{\fontsize{#1}{#1}\selectfont\color{#2}\textls[260]{FACULTY}}}
\newcommand{\acnword}[2][14]{{\fontsize{#1}{#1}\selectfont\bfseries\color{#2}accenture}}

% Integral stem with chevron crossbar. #1 centre, #2 glyph scale, #3 chevron height, #4 fg
\newcommand{\intmark}[4]{%
  \begin{scope}\gapclip{($#1+(-0.24*#3,-0.5*#3)$)}{#3}{($#1+(-0.25*#2,-0.32*#2)$)}{($#1+(0.25*#2,0.32*#2)$)}
    \node[inner sep=0,text=#4,scale=#2] at #1 {$\displaystyle\int$};
  \end{scope}
  \fill[acnPurple] \chevpath{($#1+(-0.24*#3,-0.5*#3)$)}{#3};}

% ---- the four marks, #1 = foreground colour ----
\newcommand{\markA}[1]{% 01A master lockup
  \begin{scope}\gapclip{(0.6715,0.3)}{0.9}{(0,-0.05)}{(2,1.55)}\fill[#1] \stempath;\end{scope}
  \fill[acnPurple] \chevpath{(0.6715,0.3)}{0.9};
  \node[anchor=base,inner sep=0] at (1.0,-0.83) {\facword{#1}};}
\newcommand{\markD}[1]{% 01D integral stem
  \intmark{(0,0)}{6.2}{1.0}{#1}}
\newcommand{\markE}[1]{% 01E from risk to trust
  \intmark{(0,0)}{6.2}{1.0}{#1}
  \node[anchor=west,inner sep=0,text=#1,scale=1.5] at (0.67,1.35) {$\mathit{trust}$};
  \node[anchor=west,inner sep=0,text=#1,scale=1.5] at (-0.03,-1.48) {$\mathit{risk}$};
  \node[anchor=west,inner sep=0,text=acnGrey,scale=2.2] at (0.95,-0.03) {$\mathrm{d}\mathit{AI}$};}
\newcommand{\markI}[1]{% 01I co-brand lockup
  \node[anchor=base east,inner sep=0] (a) at (-0.75,-0.22) {\acnword{#1}};
  \fill[acnPurple] \chevpath{($(a.base east)+(-0.36,0.43)$)}{0.2};
  \draw[#1,line width=0.5pt] (-0.48,-0.65) -- (-0.48,0.65);
  \intmark{(0,0)}{2.9}{0.48}{#1}
  \node[anchor=base west,inner sep=0] at (0.45,-0.22) {\facword[10]{#1}};}

\def\fgcol{white}\def\bgcol{}
\def\tmp{light}\ifx\S\tmp\def\fgcol{black}\def\bgcol{white}\fi
\def\tmp{dark}\ifx\S\tmp\def\fgcol{white}\def\bgcol{acnNight}\fi
\def\tmp{black}\ifx\S\tmp\def\fgcol{black}\fi

\begin{document}
\begin{tikzpicture}[x=14mm,y=14mm]
  \begin{scope}[local bounding box=mark]
    \csname mark\V\endcsname{\fgcol}
  \end{scope}
  \ifx\bgcol\empty
    \useasboundingbox ($(mark.south west)+(-1mm,-1mm)$) rectangle ($(mark.north east)+(1mm,1mm)$);
  \else
    \begin{scope}[on background layer]
      \fill[\bgcol] ($(mark.south west)+(-9mm,-9mm)$) rectangle ($(mark.north east)+(9mm,9mm)$);
    \end{scope}
  \fi
\end{tikzpicture}
\end{document}
"
%   \V = A (master lockup), D (integral stem), E (risk to trust), I (co-brand lockup)
%   \S = dark | light | white | black   (white/black = transparent background)
% build.sh builds every combination and exports PDF, SVG and PNG.
\documentclass[tikz,border=0pt]{standalone}
\usepackage[T1]{fontenc}
\usepackage{lmodern}
\usepackage{amsmath}
\usepackage[letterspace=120]{microtype}
\usetikzlibrary{calc,backgrounds}
\renewcommand{\familydefault}{\sfdefault}

\definecolor{acnPurple}{HTML}{A100FF}
\definecolor{acnNight}{HTML}{09090B}
\definecolor{acnGrey}{HTML}{71717A}

% House ">" (same geometry as the decks), bottom-left #1, height #2
\def\chevpath#1#2{($#1+(0,#2)$) -- ($#1+(0.62*#2,0.5*#2)$) -- #1 -- ($#1+(0,0.205*#2)$)
   -- ($#1+(0.365*#2,0.5*#2)$) -- ($#1+(0,0.795*#2)$) -- cycle}
% The same ">" grown by a 4.5% keyline (exact miter offset), clockwise
\def\chevgap#1#2{($#1+(-0.0450*#2,1.0941*#2)$) -- ($#1+(0.6917*#2,0.5000*#2)$)
   -- ($#1+(-0.0450*#2,-0.0941*#2)$) -- ($#1+(-0.0450*#2,0.2265*#2)$)
   -- ($#1+(0.2934*#2,0.5000*#2)$) -- ($#1+(-0.0450*#2,0.7735*#2)$) -- cycle}
% Clip to box #1--#2 minus the keyline area (a real gap, so transparent versions work)
\newcommand{\gapclip}[4]{\clip #3 rectangle #4 \chevgap{#1}{#2};}

% Tapered S stem (Faculty-style, redrawn). Unit box 2 x 1.5.
\def\stempath{(0,-0.018) -- (0.5,-0.03) .. controls (1.045,-0.03) and (1.045,1.47) .. (1.5,1.47)
   -- (2,1.482) -- (2,1.518) -- (1.5,1.53) .. controls (0.955,1.53) and (0.955,0.03) .. (0.5,0.03)
   -- (0,0.018) -- cycle}

\newcommand{\facword}[2][11]{{\fontsize{#1}{#1}\selectfont\color{#2}\textls[260]{FACULTY}}}
\newcommand{\acnword}[2][14]{{\fontsize{#1}{#1}\selectfont\bfseries\color{#2}accenture}}

% Integral stem with chevron crossbar. #1 centre, #2 glyph scale, #3 chevron height, #4 fg
\newcommand{\intmark}[4]{%
  \begin{scope}\gapclip{($#1+(-0.24*#3,-0.5*#3)$)}{#3}{($#1+(-0.25*#2,-0.32*#2)$)}{($#1+(0.25*#2,0.32*#2)$)}
    \node[inner sep=0,text=#4,scale=#2] at #1 {$\displaystyle\int$};
  \end{scope}
  \fill[acnPurple] \chevpath{($#1+(-0.24*#3,-0.5*#3)$)}{#3};}

% ---- the four marks, #1 = foreground colour ----
\newcommand{\markA}[1]{% 01A master lockup
  \begin{scope}\gapclip{(0.6715,0.3)}{0.9}{(0,-0.05)}{(2,1.55)}\fill[#1] \stempath;\end{scope}
  \fill[acnPurple] \chevpath{(0.6715,0.3)}{0.9};
  \node[anchor=base,inner sep=0] at (1.0,-0.83) {\facword{#1}};}
\newcommand{\markD}[1]{% 01D integral stem
  \intmark{(0,0)}{6.2}{1.0}{#1}}
\newcommand{\markE}[1]{% 01E from risk to trust
  \intmark{(0,0)}{6.2}{1.0}{#1}
  \node[anchor=west,inner sep=0,text=#1,scale=1.5] at (0.67,1.35) {$\mathit{trust}$};
  \node[anchor=west,inner sep=0,text=#1,scale=1.5] at (-0.03,-1.48) {$\mathit{risk}$};
  \node[anchor=west,inner sep=0,text=acnGrey,scale=2.2] at (0.95,-0.03) {$\mathrm{d}\mathit{AI}$};}
\newcommand{\markI}[1]{% 01I co-brand lockup
  \node[anchor=base east,inner sep=0] (a) at (-0.75,-0.22) {\acnword{#1}};
  \fill[acnPurple] \chevpath{($(a.base east)+(-0.36,0.43)$)}{0.2};
  \draw[#1,line width=0.5pt] (-0.48,-0.65) -- (-0.48,0.65);
  \intmark{(0,0)}{2.9}{0.48}{#1}
  \node[anchor=base west,inner sep=0] at (0.45,-0.22) {\facword[10]{#1}};}

\def\fgcol{white}\def\bgcol{}
\def\tmp{light}\ifx\S\tmp\def\fgcol{black}\def\bgcol{white}\fi
\def\tmp{dark}\ifx\S\tmp\def\fgcol{white}\def\bgcol{acnNight}\fi
\def\tmp{black}\ifx\S\tmp\def\fgcol{black}\fi

\begin{document}
\begin{tikzpicture}[x=14mm,y=14mm]
  \begin{scope}[local bounding box=mark]
    \csname mark\V\endcsname{\fgcol}
  \end{scope}
  \ifx\bgcol\empty
    \useasboundingbox ($(mark.south west)+(-1mm,-1mm)$) rectangle ($(mark.north east)+(1mm,1mm)$);
  \else
    \begin{scope}[on background layer]
      \fill[\bgcol] ($(mark.south west)+(-9mm,-9mm)$) rectangle ($(mark.north east)+(9mm,9mm)$);
    \end{scope}
  \fi
\end{tikzpicture}
\end{document}
